\section{Gradual type system transition from dynamic to static}

% •	Typespecs in Elixir
% •	Role of Dialyzer
% •	Limitations (lack of precision, success typing issues, etc.)
% •	Introduction to set-theoretic types

Elixir is a dynamically typed functional programming language that is currently undergoing the integration of a gradual set-theoretic type system~\cite{elixir-lang}.
Traditionally, Elixir has relied on its type specification directive, TypeSpecs, an Erlang-inspired type annotation system, to document code and support static analysis tools~\cite{elixir-typespecs}.
In particular, Dialyzer serves as the de facto standard tool for analysing TypeSpecs and detecting discrepancies through success typing~\cite{lindahl2006practical}.

However, TypeSpecs is not enforced by the compiler and therefore does not provide strong guarantees of type safety.
Instead, it primarily serves as documentation and as optional input to external analysis tools such as Dialyzer.
While Dialyzer is effective in avoiding false positives, it may produce warnings that are difficult to interpret and may lack precise information about the origin or cause of type inconsistencies.

% Compensation: TypeSpecs were designed for success typing, not sound typing.
To address these limitations, recent work on Elixir introduces a gradual set-theoretic type system known as Elixir Types~\cite{duboc2024}.
This system aims to combine the flexibility of dynamic typing with the rigor of static typing. Its benefits can be understood along three main dimensions.
First, it improves type safety guarantees by enabling more precise type inference and checking aligned with program behaviour.
Second, it supports gradual typing through the inclusion of the \texttt{dynamic()} type, allowing parts of the program to remain dynamically checked at runtime.
Third, it leverages set-theoretic constructs such as union, intersection, and negation types, enabling more expressive and compositional type definitions.

In this context, an important challenge arises: how to migrate from the existing TypeSpecs to the new Elixir Types.
To preserve semantics and maximise simplicity in the transition, this paper first investigates the translation of TypeSpecs into gradual set-theoretic types, identifying mismatches between the two systems and proposing approximation strategies where direct translation is not possible. 
Then, we bring this formulisation to implementation; a prototype is developed to evaluate the feasibility and limitations of this approach in practice.
After completing the implementation by passing test codes prescribed, we apply fine-tuning technique to train a large language model with translated types and pedict Elixir Types on untyped function definitions.